%==============================================================
%  Concurrent Model Lifecycles at the UE -- RESULTS
%
%  Revised abstract, contribution list, and populated tables.
%  Merge into the scaffold, replacing the corresponding sections.
%
%  Every number below is measured. Provenance:
%    Table I    induced.py   (nominal n=64 calib, blocked n=256)
%    Table II   experiment.py (n=512)
%    Table III  induced.py   (n=256, q_exit=25)
%    Table IV   conditional.py (n=256)
%    Table V    arbiter.py   (n=256, W_a=10)
%
%  WHAT CHANGED FROM THE SCAFFOLD, and why:
%   - The three-pipeline conflict claim did not survive. The interaction
%     that holds up is P->C. B contributes one clean secondary finding.
%   - The arbiter FAILED its own control. A blind-delay baseline beats it
%     in both bands. This is reported as a result, not omitted.
%   - The conditional check (Table IV) is what licenses that negative
%     result and is new relative to the scaffold. It belongs BEFORE the
%     arbitration section.
%==============================================================

\begin{abstract}
Release~19 makes AI/ML lifecycle management normative on the NR air
interface and organizes it per functionality: each use case receives its
own monitoring signal, decision rule, and fallback state. A user equipment,
however, runs several functionalities concurrently, driven by a common
channel realization. We construct a link-level system in which beam
management, CSI prediction, and CSI compression each carry a trained model
and a Release~19 lifecycle state machine, and we measure what
per-functionality organization does not express. Three results follow.
First, monitors calibrated independently to a common false-alarm rate fire
jointly far more often than independence predicts, and the excess grows
under degradation. Second, the cost of concurrent fallback is not additive:
along the prediction-to-compression dependency the joint cost is
\SI{26}{\percent} below the sum of the individual costs, and a prediction
fallback induces compression fallbacks that would not otherwise have fired.
Third, and contrary to our expectation, cross-functionality arbitration does
not repair this. A blind-delay baseline outperforms a dependency-aware
arbiter in both deployments, because the induced fallback action is harmful
whether or not it was induced. Measured conditional on its own monitor
firing, the compression fallback loses throughput in every degradation
regime while the prediction fallback gains. We conclude that
per-functionality lifecycle management is limited by the worst-designed
fallback in the set, and that coordination cannot substitute for fixing it.
\end{abstract}

%--------------------------------------------------------------
\subsection*{Contributions}
%--------------------------------------------------------------

\begin{itemize}
  \item A coupled three-pipeline link-level model in which beam management,
  CSI prediction, and CSI compression are driven by a single channel
  realization, each with a trained model meeting a stated acceptance
  criterion and a lifecycle state machine implementing the Release~19
  monitoring, switching, and fallback primitives
  (Sec.~\ref{sec:model}, Sec.~\ref{sec:lcm}).

  \item A measurement of \emph{correlated monitor firing} showing that
  per-functionality thresholds calibrated in isolation are jointly
  miscalibrated, that the effect is weak under nominal conditions and
  strengthens under degradation, and that its magnitude depends on how much
  of each monitor's variance is contributed by the shared channel rather
  than by its own pipeline (Sec.~\ref{sec:correlation}).

  \item An \emph{interaction} result separating genuine lifecycle coupling
  from the multiplicative composition of measurement overhead---a
  distinction that accounts for the entire apparent interaction of two of
  four subsets---and establishing sub-additive coupling along the
  prediction-to-compression dependency, together with induced fallback
  along the same path (Sec.~\ref{sec:interaction}).

  \item A \emph{negative result on arbitration}. We construct a minimal
  UE-level arbiter and a blind-delay control, and find the control
  superior in both deployments. Measured conditional on its own monitor
  firing, the compression fallback is worse than remaining active in every
  degradation regime, while the prediction fallback is better. The limiting
  factor is fallback design within a functionality, not coordination
  across functionalities (Sec.~\ref{sec:conditional},
  Sec.~\ref{sec:arbiter}).
\end{itemize}

The claim is bounded. We do not propose a replacement lifecycle framework,
we do not assert that Release~19 is defective, and the negative arbitration
result is specific to the fallback designs and operating points studied
here.

%==============================================================
\section{Results}
%==============================================================

Unless stated otherwise, FR2 is the primary deployment
(\SI{28}{\giga\hertz}, \SI{120}{\kilo\hertz} SCS, CDL-B, \SI{30}{\nano\second}
delay spread, $4\times8$ array, 128-beam oversampled codebook, Set~B of 32)
and FR1 is a second operating point (\SI{3.5}{\giga\hertz},
\SI{30}{\kilo\hertz} SCS, CDL-C, $2\times8$ array, 64 beams, Set~B of 16).
All statistics are paired at the realization level with matched seeds;
intervals are \SI{95}{\percent} percentile bootstrap.

\subsection{Model acceptance}

All three pipelines meet their acceptance criteria, fixed before the
interaction experiment and unchanged thereafter. Beam health is judged
against the exhaustive-sweep fallback it replaces rather than against a
hindsight optimum: the period-optimal beam is not attainable from a
single-slot measurement snapshot, and the non-AI path does not attain it
either.

\begin{table}[t]
\caption{Accepted model metrics (validation, 128 realizations).}
\label{tab:acceptance}
\centering
\begin{tabular}{lcc}
\toprule
 & FR2 & FR1 \\
\midrule
Beam gain loss [dB]              & 1.907  & 0.602 \\
\quad exhaustive-sweep fallback  & 1.792  & 0.435 \\
\quad excess over fallback       & 0.116  & 0.167 \\
\quad period-optimal ceiling     & 0.762  & 0.174 \\
Beam top-1                       & 0.326  & 0.604 \\
Prediction NMSE [dB]             & $-6.17$ & $-12.08$ \\
\quad hold-last NMSE [dB]        & $-1.12$ & $-5.69$ \\
\quad gain over hold-last [dB]   & 5.05   & 6.39 \\
Compression NMSE [dB]            & $-26.91$ & $-28.43$ \\
CSI report [bits]                & \multicolumn{2}{c}{64 learned / 12 fallback} \\
\bottomrule
\end{tabular}
\end{table}

\subsection{Correlated monitor firing}
\label{sec:correlation}

Each detector is calibrated in isolation to a \SI{5}{\percent} nominal
false-alarm rate, so independence predicts a joint triple rate of
$5\times10^{-2\cdot3}=1.25\times10^{-4}$.

\begin{table}[t]
\caption{Monitor firing. Ratio is measured joint rate over the product of
marginals. Nominal calibration, 64 realizations; blocked, 256.}
\label{tab:correlation}
\centering
\begin{tabular}{llcccc}
\toprule
 & & $\rho_{BP}$ & $\rho_{BC}$ & $\rho_{PC}$ & triple ratio \\
\midrule
\multirow{2}{*}{FR2} & nominal & $+0.037$ & $+0.207$ & $+0.081$ & $2.63\times$ \\
                     & blocked & $+0.210$ & $+0.236$ & $+0.023$ & $4.60\times$ \\
\multirow{2}{*}{FR1} & nominal & $+0.224$ & $+0.386$ & $+0.205$ & $20.4\times$ \\
                     & blocked & $+0.299$ & $+0.172$ & $-0.004$ & $2.06\times$ \\
\bottomrule
\end{tabular}
\end{table}

Two observations. Coincident firing exceeds the independence prediction in
every case, so per-functionality calibration is jointly optimistic. But the
magnitude is deployment-dependent rather than universal: at FR2 the
pipelines are individually noisier, which decorrelates their monitors, and
the nominal excess is modest. The FR1 blocked row must be read with care:
marginal firing there reaches \SI{25}{\percent} against a
\SI{5}{\percent} calibration target, so the detectors are saturated and
coincidence ceases to be informative---the same blockage parameters are far
more severe relative to FR1's operating point.

Note that the strongest \emph{interaction} pair, $PC$, is the weakest
\emph{correlation} pair. These are different mechanisms: $PC$ couples
through the serial dependency, not through shared observation.

\subsection{Fallback interaction}
\label{sec:interaction}

Overhead composes multiplicatively with every other cost source, so any
subset that changes measurement overhead yields a nonzero $\Delta$ by
arithmetic alone. We therefore report $\Delta$ on pre-overhead throughput.

\begin{table}[t]
\caption{Interaction $\Delta_{\text{raw}}(\mathcal{S})$, pre-overhead, 512
realizations. rel is $\Delta$ over the singleton sum. Singleton costs
$C_{\text{raw}}$: FR2 $B{=}-0.023$, $P{=}0.133$, $C{=}0.071$;
FR1 $B{=}-0.017$, $P{=}0.092$, $C{=}1.443$.}
\label{tab:interaction}
\centering
\begin{tabular}{llccc}
\toprule
 & $\mathcal{S}$ & $\Delta_{\text{raw}}$ & rel & 95\,\% CI \\
\midrule
\multirow{4}{*}{FR2}
 & $\{B,P\}$   & $+0.0090$ & $+8.2\%$  & $[+0.001,+0.018]$ \\
 & $\{B,C\}$   & $+0.0252$ & \dag      & $[+0.018,+0.033]$ \\
 & $\{C,P\}$   & $-0.0534$ & $-26.2\%$ & $[-0.059,-0.048]$ \\
 & $\{B,C,P\}$ & $-0.0189$ & \dag      & $[-0.030,-0.008]$ \\
\midrule
\multirow{4}{*}{FR1}
 & $\{B,P\}$   & $-0.0035$ & $-4.7\%$  & $[-0.010,+0.003]$ \\
 & $\{B,C\}$   & $-0.1964$ & \dag      & $[-0.218,-0.175]$ \\
 & $\{C,P\}$   & $-0.1188$ & $-7.7\%$  & $[-0.131,-0.107]$ \\
 & $\{B,C,P\}$ & $-0.3235$ & \dag      & $[-0.350,-0.298]$ \\
\bottomrule
\end{tabular}
\end{table}

\dag~Relative interaction is not reported for subsets containing $B$. The
beam fallback has \emph{negative} raw cost---exhaustive sweep is more
accurate than the model, so stripped of overhead it improves
throughput---and the singleton sum is then a near-cancellation of
opposite-signed terms, making the ratio numerically unstable.

That negative cost is itself worth stating: the beam-management fallback is
not an accuracy degradation but a resource cost. Its entire delivered cost
(FR2: 1.191 of 4.625 baseline) is the measurement overhead of sweeping 128
beams rather than 32. Per-functionality accounting, which implicitly treats
fallback as degradation, does not express this.

The result that survives in both deployments is $\{C,P\}$: sub-additive,
\SI{26}{\percent} of the predicted sum at FR2 and \SI{7.7}{\percent} at
FR1. The mechanism is saturation along the dependency chain---once
prediction has fallen back, end-to-end CSI NMSE is already 0.98, and
degrading the report further costs little.

\subsection{Induced fallback}

Matched pairs differ only in whether the source pipeline may fall back.

\begin{table}[t]
\caption{Induced fallback entries per realization, 256 realizations,
$q_{\text{exit}}{=}25$. Positive means the source \emph{induces} fallbacks
in the target; negative means it \emph{suppresses} them.}
\label{tab:induced}
\centering
\begin{tabular}{llccc}
\toprule
 & pair & induced & 95\,\% CI & $p$ \\
\midrule
\multirow{3}{*}{FR2}
 & $P\!\to\!C$ & $+0.250$ & $[+0.191,+0.309]$ & $0.0001$ \\
 & $B\!\to\!P$ & $-0.027$ & $[-0.090,+0.031]$ & $0.373$ \\
 & $B\!\to\!C$ & $+0.004$ & $[-0.043,+0.051]$ & $0.870$ \\
\midrule
\multirow{3}{*}{FR1}
 & $P\!\to\!C$ & $+0.039$ & $[+0.016,+0.062]$ & $0.0018$ \\
 & $B\!\to\!P$ & $-0.418$ & $[-0.531,-0.301]$ & $0.0001$ \\
 & $B\!\to\!C$ & $+0.070$ & $[+0.039,+0.105]$ & $0.0003$ \\
\bottomrule
\end{tabular}
\end{table}

Pairs directed into $B$ are identically zero with zero-width intervals, as
the dependency chain requires; nothing downstream feeds back into beam
selection.

$P\!\to\!C$ is the primary effect: permitting the prediction fallback
raises compression fallbacks from 0.473 to 0.723 per realization, a
\SI{53}{\percent} increase attributable to $P$'s corrective action rather
than to $C$'s own degradation. It is stable across recovery hysteresis
($+0.172$ at $q_{\text{exit}}{=}10$, $+0.250$ at 25) and weakens sixfold at
FR1, where prediction reaches \SI{-12.1}{\decibel} and its fallback
therefore damages the encoder's input far less. That scaling is the
mechanism test.

$B\!\to\!P$ at FR1 is \emph{suppression}, not induction: pinning $B$ active
increases $P$ fallbacks by 0.418. The exhaustive sweep hands the predictor
a cleaner beam. Holding a fallback is therefore not uniformly protective.

\subsection{Is the fallback the right action?}
\label{sec:conditional}

The preceding results concern when fallbacks fire and what they jointly
cost. They do not establish that firing was correct. Table~\ref{tab:conditional}
restricts the comparison to slots at which the monitor exceeded its
threshold---the regime in which a fallback is supposed to pay for itself.
Conditioning uses the statistic from the \emph{active} arm, the monitor
state that would have driven the decision; conditioning on the fallback
arm's statistic would select on an outcome the decision produced.

\begin{table}[t]
\caption{Throughput change from falling back [bit/s/Hz], positive means the
fallback helps. 256 realizations.}
\label{tab:conditional}
\centering
\begin{tabular}{llccc}
\toprule
 & & all slots & monitor fired & post-blockage \\
\midrule
\multirow{2}{*}{FR2}
 & $P$ & $-0.059$ & $+0.222$ & $+0.249$ \\
 & $C$ & $-0.072$ & $-0.269$ & $-0.258$ \\
\midrule
\multirow{2}{*}{FR1}
 & $P$ & $+0.101$ & $+0.413$ & $+0.436$ \\
 & $C$ & $-0.679$ & $-0.230$ & $-0.149$ \\
\bottomrule
\end{tabular}
\end{table}

The two pipelines separate cleanly and identically in both deployments. The
prediction fallback \emph{helps} when its monitor fires, by
\SI{0.222}{} at FR2 and \SI{0.413}{} at FR1, and helps most under deep
blockage; its episode-averaged cost is an artifact of the
\SI{93.5}{\percent} of slots on which the monitor correctly stayed quiet.
The compression fallback is \emph{worse} than remaining active whenever its
monitor fires, in every degradation phase, in both deployments. At
\SI{12}{\bit}, the wideband report is less useful than a stale learned
report even when the CSI is genuinely poor.

The specification treats these as instances of the same primitive.

\subsection{Arbitration, and why it does not help}
\label{sec:arbiter}

The arbiter withholds a downstream fallback request for at most $H$
evaluations when an upstream pipeline changed state within the last $W_a$.
The hold is bounded, and $B$ is never arbitrated, so fallback remains
locally executable and the chain cannot deadlock. As a control we run a
\emph{blind} variant that withholds every downstream request for $H$
evaluations regardless of upstream state.

\begin{table}[t]
\caption{Arbitration benefit and cost, 256 realizations, $W_a{=}10$.
``avoided'' counts induced fallbacks removed; ``denied'' counts legitimate
ones, labelled by the upstream-pinned arm.}
\label{tab:arbiter}
\centering
\begin{tabular}{llcccc}
\toprule
 & & $\Delta$ tput & avoided & denied & chain $-$ blind \\
\midrule
\multirow{2}{*}{FR2, $H{=}10$}
 & chain & $+0.020$ & 0.117 & 0.125 & \multirow{2}{*}{$-0.072$} \\
 & blind & $+0.092$ & 0.270 & 0.453 & \\
\midrule
\multirow{2}{*}{FR1, $H{=}40$}
 & chain & $+0.036$ & 0.016 & 0.043 & \multirow{2}{*}{$-0.174$} \\
 & blind & $+0.210$ & 0.039 & 0.945 & \\
\bottomrule
\end{tabular}
\end{table}

The arbiter improves throughput, and this should not be reported as
success. The blind control improves it substantially more---by
$0.072$ at FR2 and $0.174$ at FR1, intervals excluding zero
($[-0.085,-0.061]$ and $[-0.203,-0.145]$)---while \emph{denying more
legitimate fallbacks}. At $H{=}10$ the blind variant eliminates compression
fallback entirely and still wins.

When denying legitimate fallbacks improves throughput, the fallback is the
problem. Table~\ref{tab:conditional} identifies it directly: $C$'s fallback
is the wrong action whenever it is taken. Coordination cannot repair that,
because the harm does not depend on whether the trigger was induced. The
dependency-chain condition adds nothing beyond delay, and $H$ is not
binding---mean holds are 1.5--3.2 evaluations and $H\in\{10,20,40\}$ give
identical results, so the behaviour is governed by $W_a$ and the monitor
dynamics.

We report this as the outcome of a control we built to test our own
proposal. The interaction and induced-fallback results stand; the inference
from them to a need for cross-functionality arbitration does not.
